% Source: https://artofproblemsolving.com/community/c3h2419370_no_of_integers_satisfying
% Stable identity: aops-topic-2419370-q01
Move everything to one side and multiply by $2$:
\[2x^2+2y^2-2xy-4x-4y=0.\]
Adding $8$ to both sides and grouping,
\[(x-y)^2+(x-2)^2+(y-2)^2=8,\]
because $(x-y)^2+(x-2)^2+(y-2)^2=2x^2+2y^2-2xy-4x-4y+8$.

The only squares that are at most $8$ are $0$, $1$ and $4$, and the only way to write $8$ as a sum of three of them is $4+4+0$. So two of the numbers $x-y$, $x-2$, $y-2$ are $\pm2$ and the third is $0$:
\begin{itemize}
\item $x-y=0$ and $x-2=y-2=\pm2$: $(x,y)=(4,4)$ or $(0,0)$;
\item $x-2=0$, so $x=2$, and $y-2=\pm2$ (then $x-y=\mp2$ as needed): $(2,4)$ or $(2,0)$;
\item $y-2=0$, so $y=2$, and $x-2=\pm2$: $(4,2)$ or $(0,2)$.
\end{itemize}
All six pairs satisfy the equation (for example $16+4=2\cdot6+8$ for $(4,2)$), so there are $\boxed{6}$ ordered pairs.
